\documentclass[11pt,letterpaper]{article}

% =====================================================================
% THE GOODING ARCHITECTONIC STYLE GUIDE
% A Generic, Cross-Disciplinary Framework for Structural Analysis,
% Expository Writing, and Typographical Conventions.
% Applicable to the Humanities, Social Sciences, Law, Natural Sciences,
% Philosophy, and Expository Studies.
% =====================================================================

\usepackage[utf8]{inputenc}
\usepackage[T1]{fontenc}
\usepackage[english]{babel}

% --- Typography: Palatino Linotype (The Myrtlefield Standard) ---
\usepackage{mathpazo}
\usepackage{microtype}
\linespread{1.06} % Essential for the open leading required by Palatino

% --- Page Geometry ---
\usepackage[letterpaper,
            top=1in,
            bottom=1in,
            left=1in,
            right=1in,
            headheight=14pt]{geometry}

% --- Structural and Utility Packages ---
\usepackage{booktabs}
\usepackage{tabularx}
\usepackage{enumitem}
\usepackage{xcolor}
\usepackage{fancyhdr}
\usepackage{titlesec}
\usepackage{listings}
\usepackage{hyperref}

% --- Classic Myrtlefield Color Palette ---
\definecolor{myrtleWine}{RGB}{125, 22, 22}        % Deep burgundy / binding leather
\definecolor{myrtleSlate}{RGB}{70, 80, 95}        % Muted slate gray for running headers
\definecolor{myrtleDark}{RGB}{28, 28, 28}         % Traditional charcoal body ink
\definecolor{myrtleCrimson}{RGB}{145, 30, 30}     % Critical annotation / errata red
\definecolor{codeBg}{RGB}{248, 248, 246}          % Warm parchment background for code
\definecolor{codeFrame}{RGB}{220, 218, 210}       % Subtle border for listings

\hypersetup{
    colorlinks=true,
    linkcolor=myrtleWine,
    citecolor=myrtleWine,
    urlcolor=myrtleWine
}

% --- Paragraphs & Spacing ---
\setlength{\parindent}{0pt}
\setlength{\parskip}{0.45em}

% --- Running Headers & Footers ---
\pagestyle{fancy}
\fancyhf{}
\fancyhead[L]{\small\scshape The Gooding Architectonic Style Guide}
\fancyhead[R]{\small\scshape Cross-Disciplinary Analytical Method}
\fancyfoot[C]{\small\thepage}
\renewcommand{\headrulewidth}{0.4pt}

% --- Headings Formatting ---
\titleformat{\section}
  {\large\bfseries\scshape\color{myrtleDark}}
  {\thesection.}{0.6em}{}[\vspace{0.2em}\titlerule]
\titlespacing*{\section}{0pt}{18pt}{8pt}

\titleformat{\subsection}
  {\normalsize\bfseries\itshape\color{myrtleDark}}
  {\thesubsection}{0.5em}{}
\titlespacing*{\subsection}{0pt}{12pt}{4pt}

\titleformat{\subsubsection}
  {\small\bfseries\scshape\color{myrtleSlate}}
  {}{0em}{}
\titlespacing*{\subsubsection}{0pt}{8pt}{2pt}

% --- Lists Formatting ---
\setlist{leftmargin=1.8em, itemsep=0.35em, topsep=0.3em, parsep=0.15em}
\setlist[enumerate,1]{label=\textbf{\arabic*.}}
\setlist[enumerate,2]{label=(\alph*)}

% --- Code Listing Styling for LaTeX Snippets ---
\lstset{
  language=[LaTeX]TeX,
  basicstyle=\small\ttfamily\color{myrtleDark},
  backgroundcolor=\color{codeBg},
  frame=single,
  rulecolor=\color{codeFrame},
  breaklines=true,
  breakatwhitespace=true,
  tabsize=2,
  showstringspaces=false,
  keywordstyle=\bfseries\color{myrtleWine},
  commentstyle=\itshape\color{black!50}
}

% --- Editorial Macros for Scholarly Apparatus ---
\newcommand{\scholion}[2]{%
  \par\vspace{0.3em}%
  \noindent\hspace*{0.8em}%
  {\color{#1}\vrule width 1.3pt}\hspace{0.7em}%
  \parbox{\dimexpr\linewidth-2.2em}{#2}%
  \par\vspace{0.3em}%
}

\newcommand{\criticalnote}[1]{\scholion{myrtleCrimson}{\footnotesize\color{myrtleCrimson}\textbf{\textsc{Critical Note / Erratum:}} #1}}
\newcommand{\heuristicrule}[1]{\scholion{myrtleWine}{\small\color{myrtleDark}\textbf{\textsc{Core Heuristic:}} \itshape #1}}
\newcommand{\disciplinebox}[2]{\scholion{myrtleSlate}{\small\textbf{\textsc{Disciplinary Application (#1):}} #2}}

% =====================================================================
% TITLE BLOCK
% =====================================================================
\title{%
  \vspace{-1.2em}
  {\normalsize\scshape\color{myrtleSlate} General Methodology for Complex Systems and Texts}\\[0.5em]
  {\LARGE\bfseries\scshape The Gooding Architectonic Style Guide}\\[0.35em]
  {\large\itshape A Generic Framework for Structural Analysis, Expository Writing, and Typography}
}
\author{\normalsize\scshape Research \& Style Reference Manual}
\date{}

\begin{document}
\maketitle
\vspace{-1.8em}

\begin{center}
  \parbox{0.92\linewidth}{%
    \centering\itshape\small
    «The arrangement of a serious work is never an accidental ornament or a superficial aesthetic device; it is the structural scaffolding through which the human mind organizes evidence, resolves apparent contradictions, and communicates complex reality.»\\[0.3em]
    \upshape\footnotesize\scshape --- The Myrtlefield Expository Tradition
  }
\end{center}

\vspace{0.4em}
\noindent\rule{\linewidth}{0.4pt}
\vspace{0.6em}

% =====================================================================
% MANDATORY MYRTLEFIELD TABLE OF CONTENTS
% =====================================================================
\begin{center}
  {\Large\bfseries\scshape Contents}
\end{center}

\vspace{0.4em}

\begin{center}
\parbox{0.86\linewidth}{
  \small
  % --- Block 1: Front Matter (Unnumbered, Lowercase Roman Numerals) ---
  \begin{tabular*}{\linewidth}{@{\hspace{2.2em}}l@{\extracolsep{\fill}}r@{}}
    Methodological Preface & vii \\
    Taxonomy of Core Heuristics & ix \\
    List of Symmetrical Analytical Charts & xi \\
  \end{tabular*}

  \vspace{1.1em}

  % --- Block 2: Numbered Sections / Movements (Arabic Numerals) ---
  \begin{tabular*}{\linewidth}{@{}r@{\hspace{0.8em}}l@{\extracolsep{\fill}}r@{}}
    1. & Foundations of the Architectonic Method & 1 \\
    2. & The Core Analytical Heuristics across Disciplines & 3 \\
    3. & The Myrtlefield Table of Contents Specification (Mandatory) & 6 \\
    4. & The Tripartite Outline and Symmetrical Representation System & 9 \\
    5. & Expository Prose, Thought-Flow, and Critical Scholia & 13 \\
    6. & The Four-Tier Pedagogical and Review Architecture & 16 \\
    7. & Typographical and Print Design Standards & 19 \\
    8. & Generic Master \LaTeX{} Boilerplate Template & 23 \\
  \end{tabular*}

  \vspace{1.1em}

  % --- Block 3: Back Matter (Aligned with Chapter Titles) ---
  \begin{tabular*}{\linewidth}{@{\hspace{2.2em}}l@{\extracolsep{\fill}}r@{}}
    Study Guide and Methodological Practicum & 27 \\
    Cross-Disciplinary Bibliography & 30 \\
    Index of Heuristic Concepts & 33 \\
    General Analytical Index & 35 \\
  \end{tabular*}
}
\end{center}

\vspace{0.6em}
\noindent\rule{\linewidth}{0.4pt}
\vspace{1.2em}

% =====================================================================
% SECTION 1: FOUNDATIONS OF THE ARCHITECTONIC METHOD
% =====================================================================
\section{Foundations of the Architectonic Method}

\subsection{The Fallacy of Atomism}
In every intellectual discipline---whether history, law, philosophy, literature, economics, or the natural sciences---scholars and students are constantly tempted by \textbf{atomism}: treating a complex body of knowledge as a flat, disconnected catalogue of data points, definitions, or raw facts. Atomism leads inevitably to superficial memorization, conceptual vagueness, and an inability to perceive how parts integrate into a coherent whole.

\subsection{The Principle of Architectonic Coherence}
The method pioneered by Dr.~David W. Gooding (1925--2019) at Queen's University Belfast, and preserved across his philosophical, literary, and historical writings, is founded on the ancient classical insight that **form is the vehicle of meaning**. A complex text, historical epoch, or theoretical system does not consist of random occurrences; it exhibits deliberate **selection, proportion, and arrangement**.

\heuristicrule{Never present material as a disordered aggregation of isolated facts. Every individual claim, historical datum, or experimental result must be anchored within its deliberate structural movement and conceptual architecture.}

\subsection{Cross-Disciplinary Generality}
While originally developed for classical philology and biblical literature, Gooding's heuristics apply universally across diverse fields:
\begin{itemize}
  \item \textbf{History and Geopolitics:} Discerning the structural transition between the secular foundation of an empire and its enduring cultural/spiritual institutions.
  \item \textbf{Philosophy and Ethics:} Tracing the movement from foundational epistemology/metaphysics (first principles) to concrete moral and political applications.
  \item \textbf{Law and Jurisprudence:} Contrasting constitutional frameworks (statutory principles) with their symmetrical adjudication in common-law case lines.
  \item \textbf{Natural and Social Sciences:} Comparing theoretical predictive models against empirical anomalies and experimental confirmations.
  \item \textbf{Literary and Rhetorical Criticism:} Analyzing dramatic plot structures, chiastic arrangements, and character foil pairings.
\end{itemize}

% =====================================================================
% SECTION 2: THE CORE ANALYTICAL HEURISTICS ACROSS DISCIPLINES
% =====================================================================
\section{The Core Analytical Heuristics across Disciplines}

\subsection{Heuristic I: The Thought-Flow (\emph{Gedankengang}) and Boundary-Markers}
Analytical exposition must trace the dynamic movement of thought from its premises to its conclusions. This requires identifying:
\begin{enumerate}
  \item \textbf{Boundary-Markers:} Formal transition markers indicating that an author or historical process is closing one unit of thought and commencing another (e.g., temporal summaries, geographical shifts, statutory resets, or change of variables).
  \item \textbf{Paired Units and Progressions:} How consecutive incidents, cases, or lemmas are placed side-by-side to establish comparison, progression, or dialectical tension.
  \item \textbf{Internal Proportion:} The relative volume of exposition devoted to an item signals its systemic importance, independent of its chronological duration.
\end{enumerate}

\subsection{Heuristic II: Symmetrical Parallelism and Twin-Panels}
Major intellectual works and historical developments naturally tend toward balanced, two-panel or multi-panel macrostructures:
\begin{itemize}
  \item \textbf{Twin-Panel Macrostructure (Part I $\longleftrightarrow$ Part II):} A foundational panel establishing the baseline problem or secular framework, followed by a matching panel establishing its transformation, inward application, or institutional survival.
  \item \textbf{Concentric / Chiastic Inversion ($A - B - C - C' - B' - A'$):} A symmetrical structure where opening propositions recur in reverse order, placing the climactic thesis at the structural center ($C - C'$).
  \item \textbf{Side-by-Side Symmetrical Alignment ($1 \leftrightarrow 1'$, $2 \leftrightarrow 2'$):} Comparing corresponding stages point-by-point to expose thematic parallels.
\end{itemize}

\subsection{Heuristic III: The Differential Diagnosis of Recurring Strings}
When similar phenomena recur within a system (e.g., repeated economic recessions; recurring constitutional crises; parallel scientific experiments; repeated literary motifs), apply a **three-step differential diagnosis**:
\begin{enumerate}
  \item \textbf{Establish the Common Baseline:} Formulate precisely what all instances share.
  \item \textbf{Isolate the Significant Variations:} Note the minute, deliberate differences in mechanics, context, or outcome.
  \item \textbf{Deduce the Structural Function:} Explain why this specific variation appears at this precise juncture in the overarching system.
\end{enumerate}

\subsection{Heuristic IV: The Cumulative Argument and Resolution of Objections}
An architectonic exposition does not evade counterarguments; it states them in their strongest possible form. It demonstrates that superficial objections (e.g., «symmetry implies artificial distortion of reality») often rest on false dichotomies: rigorous structure and factual authenticity are fully compatible.

% =====================================================================
% SECTION 3: THE MYRTLEFIELD TABLE OF CONTENTS SPECIFICATION
% =====================================================================
\section{The Myrtlefield Table of Contents Specification (Mandatory)}

\heuristicrule{Every comprehensive document, monograph, study guide, or analytical syllabus compiled under this style guide MUST include a Table of Contents (Contents). The Table of Contents is not a mechanical page-finder, but the primary visual blueprint of the system's architecture.}

\subsection{Mandatory Structural Function}
The Table of Contents fulfills three structural purposes:
\begin{enumerate}
  \item \textbf{Macro-Symmetry at a Glance:} It displays the bipartite or multipartite symmetry of the entire work before the reader confronts individual details.
  \item \textbf{Scope Delimitation:} Every chapter or movement must carry parenthetical scope boundaries (e.g., date ranges, statutory articles, question ranges, or verse references).
  \item \textbf{Apparatus Integration:} It explicitly connects the analytical body to prefatory matter, pedagogical study guides, bibliographies, and indices.
\end{enumerate}

\subsection{The Tripartite Layout Standard}
As established in Myrtlefield House publications, the Table of Contents comprises three distinct typographical blocks:

\begin{enumerate}
  \item \textbf{Block 1: Front Matter (Preliminary Material):}
  \begin{itemize}
    \item Unnumbered entries (e.g., \emph{List of Illustrations/Tables}, \emph{List of Abbreviations}, \emph{Preface}).
    \item Indented so that the text aligns precisely with the titles of the numbered chapters below.
    \item Page numbers set in lowercase Roman numerals (`viii`, `ix`, `xi`) flush to the right margin.
  \end{itemize}

  \item \textbf{Block 2: Core Body Movements (Numbered Chapters):}
  \begin{itemize}
    \item Numbered with Arabic numerals followed by a period (`1.`, `2.`, ..., `14.`).
    \item Tabular number column ensuring identical indentation.
    \item Descriptive thematic titles with parenthetical scopes.
    \item Page numbers in Arabic numerals flush right.
    \item Separated from Block~1 and Block~3 by a distinct vertical space (`\vspace{1.1em}`).
  \end{itemize}

  \item \textbf{Block 3: Back Matter (Pedagogical and Scholarly Apparatus):}
  \begin{itemize}
    \item Unnumbered entries: \emph{Study Guide}, \emph{Appendices}, \emph{Bibliography}, and specialized \emph{Indices}.
    \item Indented to match the text column of Block~2, with Arabic page numbers flush right.
  \end{itemize}
\end{enumerate}

\subsection{Typographical Rules for the Table of Contents}
\begin{itemize}
  \item \textbf{Centered Header:} The title must be centered as \textbf{\Large Contents} (or language equivalent: \textbf{\Large Contenido}, \textbf{\Large Sommaire}).
  \item \textbf{Strict Prohibition of Dotted Leaders:} Dotted leader lines (`. . . . . 17`) are **strictly forbidden**. The Myrtlefield aesthetic demands clean, open whitespace between the title and the page number.
  \item \textbf{Typeface:} Pure Palatino Linotype with upright Roman titles and proportional numerals.
\end{itemize}

\subsection{\LaTeX{} Implementation of the Table of Contents}
The canonical tripartite layout is achieved via a dedicated `tabular*` block:

\begin{lstlisting}
\begin{center}
  {\Large\bfseries\scshape Contents}
\end{center}
\vspace{0.4em}
\begin{center}
\parbox{0.86\linewidth}{
  \small
  % Block 1: Front Matter
  \begin{tabular*}{\linewidth}{@{\hspace{2.2em}}l@{\extracolsep{\fill}}r@{}}
    Methodological Preface & vii \\
    List of Analytical Charts & ix \\
  \end{tabular*}
  \vspace{1.1em}
  % Block 2: Numbered Movements
  \begin{tabular*}{\linewidth}{@{}r@{\hspace{0.8em}}l@{\extracolsep{\fill}}r@{}}
    1. & Foundational Theory and Axioms (Scope A) & 1 \\
    2. & Concrete Adjudication and Cases (Scope B) & 18 \\
    3. & Empirical Verification and Synthesis (Scope C) & 35 \\
  \end{tabular*}
  \vspace{1.1em}
  % Block 3: Back Matter
  \begin{tabular*}{\linewidth}{@{\hspace{2.2em}}l@{\extracolsep{\fill}}r@{}}
    Study Guide & 52 \\
    Bibliography & 58 \\
    General Analytical Index & 62 \\
  \end{tabular*}
}
\end{center}
\end{lstlisting}

% =====================================================================
% SECTION 4: THE TRIPARTITE OUTLINE AND SYMMETRICAL REPRESENTATION SYSTEM
% =====================================================================
\section{The Tripartite Outline and Symmetrical Representation System}

A complete Gooding outline consists of three levels of analytical representation:

\subsection{Level 1: The General Plan (\emph{Plan General})}
The macro-level summary dividing the entire topic into balanced Parts:

\begin{center}
\small
\begin{tabularx}{\linewidth}{@{}p{0.48\linewidth}@{\hspace{0.04\linewidth}}p{0.48\linewidth}@{}}
  \toprule
  \multicolumn{1}{c}{\textbf{\textsc{Panel I: Foundational Framework}}} & 
  \multicolumn{1}{c}{\textbf{\textsc{Panel II: Applied Transformation}}} \\
  \multicolumn{1}{c}{\footnotesize\color{myrtleSlate}(Theory, institutional setup, initial conflict)} & 
  \multicolumn{1}{c}{\footnotesize\color{myrtleSlate}(Practice, internal maturation, enduring survival)} \\
  \midrule
  \textbf{Stage 1: Origin and Structural Premises}\newline
  The external confrontation, dismantling of prior atomism, and establishment of first principles. &
  \textbf{Stage 3: Concrete Implementation and Culture}\newline
  The internal consolidation: institutional mechanisms, ethical models, and behavioral paradigms. \\
  \addlinespace[0.5em]
  \textbf{Stage 2: Institutional Development}\newline
  Conflict between competing systems, legal codification, and territorial/administrative ordering. &
  \textbf{Stage 4: Culmination and Systemic Survival}\newline
  Resolution of core tensions, survival through systemic crisis, and legacy for subsequent generations. \\
  \bottomrule
\end{tabularx}
\end{center}

\subsection{Level 2: The Symmetrical Structural Table (\emph{Tabla Estructural})}
The meso-level comparative table formatted strictly according to **booktabs** principles:
\begin{itemize}
  \item Zero vertical rules.
  \item Precise cross-column alignment between corresponding stages ($A \longleftrightarrow A'$).
  \item Parallel alphanumeric labelling across panels.
\end{itemize}

\begin{center}
\footnotesize
\begin{tabularx}{\linewidth}{@{}l X l X@{}}
  \toprule
  \multicolumn{2}{c}{\textbf{\textsc{Phase A: Theoretical Formulation}}} & 
  \multicolumn{2}{c}{\textbf{\textsc{Phase B: Empirical Application}}} \\
  \midrule
  \textbf{1.} & \textbf{Initial Thesis:} Formulation of basic axioms; definition of foundational boundaries. &
  \textbf{1$'$.} & \textbf{Final Synthesis:} Experimental verification; enduring model surviving boundary breakdown. \\
  \addlinespace[0.4em]
  \textbf{2.} & \textbf{Systemic Friction:} Resistance from competing paradigms or hostile legal regimes. &
  \textbf{2$'$.} & \textbf{Paradigm Integration:} Assimilation of competing mechanisms into a unified framework. \\
  \bottomrule
\end{tabularx}
\end{center}

\subsection{Level 3: The Detailed Analytical Outline (\emph{Esquema Analítico})}
The micro-level hierarchical breakdown. Every entry must:
\begin{enumerate}
  \item Provide a bold thematic title with bracketed question/section scopes.
  \item Articulate the **causal nexus** (the *why* behind the movement of thought).
  \item Explicitly flag critical discrepancies, counter-theses, or examination warnings.
\end{enumerate}

% =====================================================================
% SECTION 5: EXPOSITORY PROSE, THOUGHT-FLOW, AND CRITICAL SCHOLIA
% =====================================================================
\section{Expository Prose, Thought-Flow, and Critical Scholia}

\subsection{The Expository Register}
Expository prose in the Gooding tradition balances **scientific precision** with **pedagogical lucidity**:
\begin{enumerate}
  \item \textbf{Precision of Nomenclature:} Technical terms must never be conflated (e.g., distinguishing civil statutes from common-law maxims; distinguishing an intermediate legal privilege from full citizenship).
  \item \textbf{Logical Connectives:} Use measured deductive phrases: «Judged merely from the formal arrangement...», «This observation calls for comment...», «The conclusion is inescapable that...».
\end{enumerate}

\subsection{The Categorization of Scholarly Annotations}
Annotations must be visually and functionally segregated:
\begin{itemize}
  \item \textbf{Provenance Citations (`\textbackslash fuente`):} Flush-right bracketed references in small caps: `[\textsc{source}~4; \textsc{ref}~12]`.
  \item \textbf{Critical Scholia and Errata (`\textbackslash criticalnote` / `\textbackslash examen`):} Set as an indented paragraph with a left vertical border in deep crimson (`myrtleCrimson`). Used for correcting misconceptions, past exam errors, or fallacious interpretations.
  \item \textbf{Discrepancy and Textual Notices (`\textbackslash aviso`):} Set with a neutral slate vertical bar, highlighting tensions between primary documents and secondary commentary.
\end{itemize}

% =====================================================================
% SECTION 6: THE FOUR-TIER PEDAGOGICAL AND REVIEW ARCHITECTURE
% =====================================================================
\section{The Four-Tier Pedagogical and Review Architecture}

Review sections and study guides must not present random question pools. They must structure inquiry into four cognitive tiers:

\begin{enumerate}[label=\textbf{\Roman*.}]
  \item \textbf{Tier I: Fundamental Definitions and Conceptual Taxonomy:} Precise definition of foundational terms, opposing pairs (e.g., *Integration* vs. *Atomism*), and technical categories.
  \item \textbf{Tier II: Specific Institutional and Causal Queries:} Rigorous cause-and-effect questions tracing historical chronology, mechanical steps, or statutory sequences with cross-references `[§X]`.
  \item \textbf{Tier III: Formal Analysis of Cases, Artifacts, or Experiments:} A four-parameter diagnostic template:
  \begin{itemize}
    \item \textbf{Formal Designation:} Official name, title, or citation.
    \item \textbf{Geographical / Systemic Locus:} Context, jurisdiction, or apparatus.
    \item \textbf{Chronology / Epoch:} Date, phase, or sequence number.
    \item \textbf{Distinctive Formal Feature:} The unique trait, holding, or finding.
  \end{itemize}
  \item \textbf{Tier IV: Integrative Synthesis and Architectural Essays:} Broad prompts requiring the analyst to reconstruct the entire thought-flow from memory without auxiliary notes.
\end{enumerate}

% =====================================================================
% SECTION 7: TYPOGRAPHICAL AND PRINT DESIGN STANDARDS
% =====================================================================
\section{Typographical and Print Design Standards}

\subsection{Typeface: Palatino Linotype}
Palatino is the mandatory typeface of the Myrtlefield canon. It provides robust serifs, legible proportions, and historical warmth. In \LaTeX{}:
\begin{lstlisting}
\usepackage{mathpazo}
\usepackage{microtype}
\linespread{1.06} % Compensates for Palatino's generous x-height
\end{lstlisting}

\subsection{Geometry and Margin Proportions}
Documents must maintain dignified, book-like proportions:
\begin{lstlisting}
\usepackage[letterpaper, top=1in, bottom=1in, left=1in, right=1in, headheight=14pt]{geometry}
\end{lstlisting}

\subsection{Running Headers and Footers}
Running headers reinforce the macrostructure:
\begin{lstlisting}
\usepackage{fancyhdr}
\pagestyle{fancy}
\fancyhf{}
\fancyhead[L]{\small\scshape Overall Topic / Course / Book Title}
\fancyhead[R]{\small\scshape Current Chapter / Movement Title}
\fancyfoot[C]{\small\thepage}
\renewcommand{\headrulewidth}{0.4pt}
\end{lstlisting}

\subsection{Macro Library for Academic Apparatus}
\begin{lstlisting}
\definecolor{myrtleWine}{RGB}{125, 22, 22}
\definecolor{myrtleCrimson}{RGB}{145, 30, 30}
\definecolor{myrtleSlate}{RGB}{70, 80, 95}

% Source Citation
\newcommand{\fuente}[1]{\nolinebreak\hfill{\footnotesize\color{black!55}\textsc{[#1]}}}

% Question / Lemma Item
\newcommand{\preg}[1]{\item \textbf{#1}\par\nopagebreak\vspace{0.2em}}

% Critical Scholion / Erratum
\newcommand{\criticalnote}[1]{%
  \par\vspace{0.35em}%
  \noindent\hspace*{0.8em}%
  {\color{myrtleCrimson}\vrule width 1.3pt}\hspace{0.7em}%
  \parbox{\dimexpr\linewidth-2.2em}{%
    \footnotesize\color{myrtleCrimson}%
    \textbf{\textsc{Critical Note / Erratum:}} #1%
  }%
  \par\vspace{0.35em}%
}

% Discrepancy Notice
\newcommand{\aviso}[1]{%
  \par\vspace{0.3em}%
  \noindent\hspace*{0.8em}%
  {\color{myrtleSlate}\vrule width 1.1pt}\hspace{0.7em}%
  \parbox{\dimexpr\linewidth-2.2em}{%
    \footnotesize\color{black!75}%
    \textbf{\textsc{Textual / Methodological Notice:}} #1%
  }%
  \par\vspace{0.3em}%
}
\end{lstlisting}

% =====================================================================
% SECTION 8: GENERIC MASTER LATEX BOILERPLATE TEMPLATE
% =====================================================================
\section{Generic Master \LaTeX{} Boilerplate Template}

The following standalone boilerplate is applicable to any academic discipline:

\begin{lstlisting}
\documentclass[11pt,letterpaper]{article}

\usepackage[utf8]{inputenc}
\usepackage[T1]{fontenc}
\usepackage[english]{babel}
\usepackage{mathpazo}
\usepackage{microtype}
\linespread{1.06}

\usepackage[letterpaper, margin=1in, headheight=14pt]{geometry}
\usepackage{booktabs, tabularx, enumitem, xcolor, fancyhdr, titlesec, hyperref}

\definecolor{myrtleWine}{RGB}{125, 22, 22}
\definecolor{myrtleCrimson}{RGB}{145, 30, 30}
\definecolor{myrtleSlate}{RGB}{70, 80, 95}

\pagestyle{fancy}
\fancyhf{}
\fancyhead[L]{\small\scshape Discipline / Research Project}
\fancyhead[R]{\small\scshape Chapter Title}
\fancyfoot[C]{\small\thepage}
\renewcommand{\headrulewidth}{0.4pt}

\titleformat{\section}{\large\bfseries\scshape}{\thesection.}{0.6em}{}[\vspace{0.2em}\titlerule]
\setlist[enumerate,1]{label=\textbf{\arabic*.}, leftmargin=1.8em, itemsep=0.8em}

\newcommand{\fuente}[1]{\nolinebreak\hfill{\footnotesize\color{black!55}\textsc{[#1]}}}
\newcommand{\preg}[1]{\item \textbf{#1}\par\nopagebreak\vspace{0.2em}}

\title{%
  {\normalsize\scshape Department / Research Division}\\[0.4em]
  {\LARGE\bfseries\scshape Architectural Study of the System}\\[0.3em]
  {\large\itshape An Expository Analysis and Structural Guide}
}
\author{\normalsize\scshape Researcher / Author Name}
\date{}

\begin{document}
\maketitle

% 1. Epigraph / Overarching Thesis
\begin{center}
  \parbox{0.88\linewidth}{\centering\itshape\small «Central research question or thesis governing the system.»}
\end{center}

% 2. MANDATORY MYRTLEFIELD TABLE OF CONTENTS
\begin{center}
  {\Large\bfseries\scshape Contents}
\end{center}
\vspace{0.3em}
\begin{center}
\parbox{0.86\linewidth}{
  \small
  % Block 1: Front Matter
  \begin{tabular*}{\linewidth}{@{\hspace{2.2em}}l@{\extracolsep{\fill}}r@{}}
    Methodological Framework and Sources & 1 \\
    General Plan and Structural Symmetry & 2 \\
  \end{tabular*}
  \vspace{0.8em}
  % Block 2: Core Movements
  \begin{tabular*}{\linewidth}{@{}r@{\hspace{0.8em}}l@{\extracolsep{\fill}}r@{}}
    1. & Foundational Architecture (Scope A) & 3 \\
    2. & Concrete Mechanisms and Applications (Scope B) & 7 \\
  \end{tabular*}
  \vspace{0.8em}
  % Block 3: Back Matter
  \begin{tabular*}{\linewidth}{@{\hspace{2.2em}}l@{\extracolsep{\fill}}r@{}}
    Study Guide and Methodological Review & 12 \\
    Bibliographical Apparatus & 15 \\
  \end{tabular*}
}
\end{center}

\vspace{0.4em}
\noindent\rule{\linewidth}{0.4pt}

% 3. General Plan and Symmetrical Table
\section*{General Plan and Symmetrical Analysis}
\begin{center}
\small
\begin{tabularx}{\linewidth}{@{}p{0.48\linewidth}@{\hspace{0.04\linewidth}}p{0.48\linewidth}@{}}
  \toprule
  \multicolumn{1}{c}{\textbf{\textsc{Panel I: Theoretical Foundation}}} & 
  \multicolumn{1}{c}{\textbf{\textsc{Panel II: Applied Transformation}}} \\
  \midrule
  \textbf{Stage 1: Axiomatic Formulation} & \textbf{Stage 3: Applied Execution} \\
  \textbf{Stage 2: Institutional Friction} & \textbf{Stage 4: Enduring Synthesis} \\
  \bottomrule
\end{tabularx}
\end{center}

% 4. Expository Body
\section{Foundational Architecture}
\begin{enumerate}
  \preg{What is the core axiom governing the system?}
  The foundational principle is that form reflects function... \fuente{Primary~Source}
\end{enumerate}

\end{document}
\end{lstlisting}

\vfill
\noindent\rule{\linewidth}{0.4pt}\\[0.2em]
{\footnotesize\color{black!60}%
\textbf{\textsc{Colophon:}} Formulated by distilling the literary, philosophical, and structural heuristics of Dr.~David W. Gooding into a generic, cross-disciplinary methodology for academic research, legal analysis, scientific exposition, and textual hermeneutics. Set in \LaTeX{} with Palatino Linotype.}

\end{document}
